% =====================================================================
% Wave-13 B1: Frontier automorphic forms from the last 10 slides of
% Lorgat 2020, beyond the 8 Gritsenko-Clery paramodular diagonal-divisor
% forms. Reads slides 55-64 (lines 800-953 of wave13_source_slides.txt)
% against the paper's Conjecture 1 (paper p.1, lines 108-121 of
% wave13_source_paper.txt) and the Arithmetic Geometry of CY3 section.
%
% Mission. The 8 GC forms of slides 47-54 are classification data:
% paramodular forms vanishing to order 1 along diagonal translates.
% The frontier is the Borcherds-lift image: each of the 8 arises as
% the reciprocal-square-root of a twisted DT zeta function of a CY3
% quotient X=(S x E)/(Z/NZ), with Borcherds-input the (g_N,h_M)-twisted
% K3 elliptic genus F^{(g_N,h_M)}_L. The last 10 slides name the
% case enumeration N in {1,2,3,5,7} / N=4 / N=6 / N=8 that drives
% the frontier superalgebras g_{L,(g_N,h_M)} extending g_{Delta_5}.
%
% Five ATTACK<->HEAL cycles. Chriss-Ginzburg voice. Honest \ClaimStatus*.
% Cross-references the two-stage functor Phi = Sp_{Sigma_{d-1},C} o Phi^{FA}_d
% (notes/wave12_u1_two_stage_functor.tex) and kappa_BKM(Phi_N)=c_N(0)/2
% universality (notes/wave12_a2_k3xE_BKM_kazhdan.tex).
%
% Authorship: Raeez Lorgat.
% =====================================================================

\providecommand{\kBKM}{\kappa_{\mathrm{BKM}}}
\providecommand{\kch}{\kappa_{\mathrm{ch}}}
\providecommand{\kcat}{\kappa_{\mathrm{cat}}}
\providecommand{\kfib}{\kappa_{\mathrm{fiber}}}
\providecommand{\ClaimStatusTheorem}{\textbf{[Thm]}}
\providecommand{\ClaimStatusDerivation}{\textbf{[derivable]}}
\providecommand{\ClaimStatusConjectured}{\textbf{[Conj]}}
\providecommand{\ClaimStatusDefinition}{\textbf{[Def]}}
\providecommand{\ClaimStatusOpen}{\textbf{[Open]}}
\providecommand{\Sp}{\mathrm{Sp}}
\providecommand{\OO}{\mathrm{O}}
\providecommand{\Hb}{\mathbb{H}}
\providecommand{\fg}{\mathfrak{g}}
\providecommand{\fh}{\mathfrak{h}}
\providecommand{\fn}{\mathfrak{n}}
\providecommand{\cO}{\mathcal{O}}
\providecommand{\CoHA}{\mathrm{CoHA}}
\providecommand{\mult}{\mathrm{mult}}
\providecommand{\Hilb}{\mathrm{Hilb}}
\providecommand{\DT}{\mathrm{DT}}
\providecommand{\Jac}{\mathrm{Jac}}
\providecommand{\LII}{\Lambda^{2,1}_{I\!I}}
\providecommand{\Lthrtwo}{\Lambda^{3,2}}
\providecommand{\Weyl}{W}
\providecommand{\Aut}{\mathrm{Aut}}
\providecommand{\Nkl}[1]{N\!=\!#1}

\section*{Wave 13 B1: Frontier automorphic forms of Lorgat 2020
beyond the eight Gritsenko--Cl\'ery paramodular diagonal-divisor forms}
\label{sec:wave13-b1-frontier-forms}

The last decade of slides (55--64) unfolds a single construction and a
single conjecture.
Given a K3 surface $S$ equipped with a pair of finite-order symplectic
automorphisms $(g_N,h_M)$ with $N,M\leq 8$ (Nikulin bound \cite{Nikulin84}),
a lattice polarisation $L\subset\mathrm{NS}(S)$ fixed by both, and an
elliptic curve $(E,e_0)$ with an $N$-torsion point,
the quotient $X=(S\times E)/(\mathbb{Z}/N\mathbb{Z})$ by the free
diagonal action $(s,e)\mapsto(g_N s,e+e_0)$ is a smooth projective
Calabi--Yau threefold.
A curve-class $\beta_h=\tfrac1N(s_1+\cdots+s_N+hF)\in\mathrm{Pic}(X)$
organises the reduced Donaldson--Thomas generating series
\[
Z^X_{L,h_M}(q,t,p)
\;=\;
\sum_{h\geq0}\sum_{d\geq0}\sum_{n\in\mathbb Z}
\DT^X_{n,(\beta_h,d)}\,
q^{d-1+\tfrac12\langle\beta_h,\beta_h\rangle}\,
t^d\,(-p)^n,
\]
twisted by the $h_M$-cover class.
The eight Gritsenko--Cl\'ery diagonal-divisor forms
\[
\bigl\{\Delta_5,\;M_2(\Gamma_2,\nu_4),\;M_3(\Gamma_1(2),\nu_2),\;
M_1(\Gamma_3,\nu_6),\;M_2(\Gamma_1(3),\nu_2),\;
M_{1/2}(\Gamma_4,\nu_8),\;M_{3/2}(\Gamma_1(4),\nu_4),\;
M_1(\Gamma_2(2),\nu_4)\bigr\}
\]
are claimed by \emph{Conjecture~1} to arise, up to constant, as
$(-Z^X_{L,h_M})^{-1/2}$ for eight matched data $(S,L,g_N,h_M,E,e_0)$,
each denominator-function of a frontier BKM superalgebra
$\fg_{L,(g_N,h_M)}$ with root multiplicities read off the
$(g_N,h_M)$-twisted K3 elliptic genus $F^{(g_N,h_M)}_L$, a weight-$0$,
index-$1$ Jacobi form.
The ``frontier'' is the seven sibling Borcherds lifts beyond
$N=1\to\Delta_5$, indexed by the Nikulin orbifold spectrum
$N\in\{2,3,4,5,6,7,8\}$, that the slides announce case-by-case and
leave open at the level of denominator verification against the
Oberdieck--Pixton DT zeta function.

\subsection*{The frontier list}
\label{sec:frontier-list}

\begin{table}[h]
\small\centering
\begin{tabular}{c|l|l|l|l|l|l}
$N$ & Frontier form & Weight & Level / group & Borcherds input
$F^{(g_N,h_M)}_L$ & $\fg_L\subset\fg_{\Delta_5}$ & Claim status\\
\hline
$1$ & $\Delta_5$ & $5$ & $\Sp_4(\mathbb Z)$, $\nu_2$
& $F^{(\mathrm{id},\mathrm{id})}=\phi_{0,1}$
& $\fg_{\Delta_5}$ (Lorgat 2020 Thm.~3)
& \ClaimStatusTheorem\\
$2$ & $M_2(\Gamma_2,\nu_4)$ & $2$ & paramodular $\Gamma_2$, $\nu_4$
& $F^{(g_2,g_2)}_{L_2}$ (M$_{24}$-frame shape $1^8 2^8$)
& $\fg_{L_2,(g_2,g_2)}\subset\fg_{\Delta_5}$
& \ClaimStatusConjectured\\
$2'$ & $M_3(\Gamma_1(2),\nu_2)$ & $3$ & Hecke $\Gamma_1(2)$, $\nu_2$
& $F^{(g_2,\mathrm{id})}_{L_2}$ (pair with non-split pair)
& $\fg_{L_2,(g_2,\mathrm{id})}\subset\fg_{\Delta_5}$
& \ClaimStatusConjectured\\
$3$ & $M_1(\Gamma_3,\nu_6)$ & $1$ & paramodular $\Gamma_3$, $\nu_6$
& $F^{(g_3,g_3)}_{L_3}$ (M$_{24}$-frame shape $1^6 3^6$)
& $\fg_{L_3,(g_3,g_3)}\subset\fg_{\Delta_5}$
& \ClaimStatusConjectured\\
$3'$ & $M_2(\Gamma_1(3),\nu_2)$ & $2$ & Hecke $\Gamma_1(3)$, $\nu_2$
& $F^{(g_3,\mathrm{id})}_{L_3}$
& $\fg_{L_3,(g_3,\mathrm{id})}\subset\fg_{\Delta_5}$
& \ClaimStatusConjectured\\
$4$ & $M_{1/2}(\Gamma_4,\nu_8)$ & $1/2$ & paramodular $\Gamma_4$, $\nu_8$
& $F^{(g_4,g_4)}_{L_4}$ (M$_{24}$-frame shape $2^4 4^4$)
& $\fg_{L_4,(g_4,g_4)}\subset\fg_{\Delta_5}$
& \ClaimStatusConjectured\\
$4'$ & $M_{3/2}(\Gamma_1(4),\nu_4)$ & $3/2$ & Hecke $\Gamma_1(4)$, $\nu_4$
& $F^{(g_4,\mathrm{id})}_{L_4}$
& $\fg_{L_4,(g_4,\mathrm{id})}\subset\fg_{\Delta_5}$
& \ClaimStatusConjectured\\
$2''$ & $M_1(\Gamma_2(2),\nu_4)$ & $1$ & Hecke-in-paramodular $\Gamma_2(2)$, $\nu_4$
& $F^{(g_2,h_2)}_{L_2}$ (split pair, $L_2$ of rank $16$)
& $\fg_{L_2,(g_2,h_2)}\subset\fg_{\Delta_5}$
& \ClaimStatusConjectured\\
\hline
$5$ & $M_{?}(\Gamma_5,\nu_?)$ & $c_5(0)/2$ & $\Gamma_5$
& $F^{(g_5,g_5)}_{L_5}$ (frame shape $1^4 5^4$)
& $\fg_{L_5,(g_5,g_5)}$
& \ClaimStatusOpen\\
$6$ & $M_{?}(\Gamma_6,\nu_?)$ & $c_6(0)/2$ & $\Gamma_6$
& $F^{(g_6,g_6)}_{L_6}$ (frame shape $1\,2\,3\,6$)
& $\fg_{L_6,(g_6,g_6)}$
& \ClaimStatusOpen\\
$7$ & $M_{?}(\Gamma_7,\nu_?)$ & $c_7(0)/2$ & $\Gamma_7$
& $F^{(g_7,g_7)}_{L_7}$ (frame shape $1^3 7^3$)
& $\fg_{L_7,(g_7,g_7)}$
& \ClaimStatusOpen\\
$8$ & $M_{?}(\Gamma_8,\nu_?)$ & $c_8(0)/2$ & $\Gamma_8$
& $F^{(g_8,g_8)}_{L_8}$ (frame shape $1\,2\,4\,8$)
& $\fg_{L_8,(g_8,g_8)}$
& \ClaimStatusOpen\\
\end{tabular}
\caption{Frontier automorphic forms: the eight Gritsenko--Cl\'ery
diagonal-divisor forms identified on slides 47--54 (cases
$N\in\{1,2,3,4\}$ and their Hecke or split-pair variants), and four
open frontier extensions at $N\in\{5,6,7,8\}$ not realised by the GC
classification but predicted as denominator-functions of siblings
$\fg_{L_N,(g_N,g_N)}\subset\fg_{\Delta_5}$ via the
$(g_N,h_M)$-twisted-elliptic-genus Borcherds lift.
The weight column records the Siegel/paramodular weight of the form;
the Borcherds weight $\kBKM=c_N(0)/2$ identification (Borcherds 1998
Thm.~13.3) makes the weight a direct readout of the constant Fourier
coefficient of the twisted Jacobi input.}
\label{tab:frontier-forms}
\end{table}

\subsection*{Cycle 1: The $N=1$ anchor and the weight-level arithmetic}
\label{sec:cycle-1}

\paragraph{Attack.}
The slide reads $\Delta_5\in M_5(\Gamma_1,\nu_2)$; the paper writes
$\Delta_5\in M_5(\Sp_4(\mathbb Z),\nu_{\Delta_5})$ with the multiplier
system $\nu_{\Delta_5}$ of order 2.
Which group, and in what sense is the weight $5$?
Three potential confusions appear:
the Siegel weight (form on $\Sp_4(\mathbb Z)$-invariants of $\Hb_2$),
the paramodular weight (form on $\Gamma_t\subset\Sp_4(\mathbb Q)$),
and the Borcherds-lift weight (constant term of the Jacobi input).

\paragraph{Ghost theorem.}
The three weights coincide at $N=1$ by the isomorphism
$\Gamma_1=\Gamma_1(1)=\Sp_4(\mathbb Z)$ (paper p.~1 line 46), together
with the Maa\ss--Sp$_4$ character table \cite[\S5]{Maass1964}: the
multiplier system $\nu_2$ on $\Gamma_1$ is the unique non-trivial
order-$2$ character of $\Sp_4(\mathbb Z)$, and the form $\Delta_5$
representing its $5$-dimensional weight-space is the Maa\ss--Gritsenko
lift of $\phi_{0,1}$.

\begin{theorem}[Weight-level reconciliation at $N=1$]
\label{thm:n1-weight-reconciliation}
The Igusa cusp form $\Delta_5$ admits three equivalent presentations,
all with $\kBKM=5$:
\begin{enumerate}[label={\rm(\roman*)}]
\item as a Siegel modular form on $\Sp_4(\mathbb Z)=\Gamma_1$ of
weight $5$ with multiplier $\nu_{\Delta_5}=\nu_2$ of order $2$
\emph{(slide 55; Igusa 1964)};
\item as a paramodular form on the level-$1$ paramodular group
$\Gamma_1\subset\Sp_4(\mathbb Q)$ of weight $5$ with multiplier
$\nu_2$ \emph{(paper p.~1 Conjecture~1; Gritsenko--Nikulin 1995)};
\item as the Borcherds-singular-theta lift of the weight-$0$
index-$1$ weak Jacobi form $\phi_{0,1}$ on the lattice
$\Lthrtwo\simeq\Lambda^{1,1}\oplus\Lambda^{1,1}\oplus[2]$, of weight
$\mathrm{wt}_{\mathrm{Borch}}=c_1(0,0)/2=10/2=5$
\emph{(paper Theorem~4; Borcherds 1998 Thm.~13.3)}.
\end{enumerate}
The identity $\tfrac1{64}\Delta_5(2Z)=\Phi_{\fg_{\Delta_5}}(z)$
realises $\Delta_5$ as the automorphic denominator of the BKM
superalgebra $\fg_{\Delta_5}$ (paper Theorem~3).
\end{theorem}

\paragraph{Heal.}
The three weights agree at $N=1$ because the three groups coincide;
at $N\geq2$ they may disagree by the index factor of the paramodular
embedding $\Gamma_N\hookrightarrow\Sp_4(\mathbb Q)$, and the frontier
weights $\tfrac12,1,\tfrac32,2,3$ appear precisely for the paramodular
and Hecke presentations of the seven GC siblings.
Recording the weight as $\kBKM=c_N(0)/2$ with $c_N(0):=c_{F^{(g_N,h_M)}_L}(0,0)$
unifies the three readings.

\subsection*{Cycle 2: Twisted K3 elliptic genera as Borcherds input}
\label{sec:cycle-2}

\paragraph{Attack.}
Slide 59 claims \emph{``there is a GBKM algebra with correction
$\fg\subset\fg_L$, depending on $L$, given by $(g_N,h_M)$-twisted
elliptic genus of a K3-surface $F^{(g_N,h_M)}_L$, a Jacobi form of
weight $0$, index $1$''}.
Two points require first-principles verification:
(a) that the twisted elliptic genus is a weight-$0$ index-$1$
Jacobi form for every Nikulin-admissible $(g_N,h_M)$ pair; and
(b) that its Borcherds lift lies in the Gritsenko--Cl\'ery
diagonal-divisor table.

\paragraph{Ghost theorem.}
The Mathieu-moonshine correspondence of
Eguchi--Ooguri--Tachikawa \cite{EOT2011} identifies the
twisted K3 elliptic genus with the $M_{24}$-class character valued
at the Nikulin symplectic class $[g_N]$: for every class
$g\in M_{24}$ with Frame shape $\prod_i k_i^{a_i}$ one writes
\[
F^{(g,\mathrm{id})}(\tau,z)
=\tfrac{1}{|\mathrm{stab}(g)|}
\sum_{h\in\langle g\rangle}
\mathrm{tr}_{V^{g,h}}(-1)^{F_L+F_R}q^{L_0}\bar q^{\bar L_0}y^{J_0},
\]
and this is a weak Jacobi form of weight~$0$ index~$1$ with a
congruence character depending on the Frame shape.
The vector-valued Borcherds lift of Borcherds 1998 Theorem~13.3,
applied to the Eichler--Zagier theta decomposition of
$F^{(g_N,h_M)}_L$ on the lattice $L\oplus\Lambda^{1,1}\oplus[2]$,
is an automorphic form of weight $c_{F^{(g_N,h_M)}_L}(0,0)/2$ on
$\OO(L\oplus\Lambda^{1,1}\oplus[2])_+$, which upon restriction to
the paramodular cusp specialises to the paramodular form of level $N$
(or the Hecke-congruence form at level $\Gamma_1(N)$).

\begin{proposition}[Borcherds-input normalisation]
\label{prop:borcherds-input-normalisation}
For the Nikulin-admissible orbifold data $(g_N,h_M)$ with $N,M\leq 8$
and any $L$-polarisation compatible with $\langle g_N,h_M\rangle$, the
$(g_N,h_M)$-twisted K3 elliptic genus $F^{(g_N,h_M)}_L(\tau,z)$ is a
weak Jacobi form of weight $0$, index $1$, and its scalar-valued
$\ell=0$ constant term at $q^0$ satisfies
\[
c_N(0,0)\;=\;\tfrac{24}{|\langle g_N,h_M\rangle|}
\cdot\chi(K3)\big|_{\langle g_N,h_M\rangle\text{-fixed}}
\;+\;\text{(Hecke correction)}\;\in\;2\,\mathbb Z.
\]
The Borcherds lift $\Phi_{F^{(g_N,h_M)}_L}$ is a Siegel paramodular
form of weight $c_N(0,0)/2$ on $\Gamma_N$ with character determined
by the Frame shape.
\end{proposition}

\begin{proof}[Proof sketch]
Weak Jacobi-form type: follows from modular covariance of the
twining orbifold partition function
\cite[\S4]{Gaberdiel-Hohenegger-Volpato2012}.
Integrality of $c_N(0,0)$: follows from the Fourier-coefficient
integrality of the frame-shape character and the $\phi_{0,1}$-basis
expansion $F^{(g_N,h_M)}_L=a\phi_{0,1}+b\phi_{-2,1}\cdot E_2$ for
rational $a,b$ (weak-Jacobi ring presentation
\cite{EichlerZagier1985}).
Parity of $c_N(0,0)$: follows from $c_1(0,0)=10$ at $N=1$ and the
divisor character of the Borcherds-Ray super-correction
\cite[\S3]{Ray2007}; the character forces $c_N(0,0)\in 2\mathbb Z$
for all orbifold classes.
Borcherds-lift assembly: Borcherds 1998 Thm.~13.3 applied verbatim
with the vector-valued $\vartheta$-decomposition of
$F^{(g_N,h_M)}_L$; output weight $c_N(0,0)/2$ by direct substitution.
\end{proof}

\paragraph{Heal.}
The universal Borcherds weight identity
$\kBKM(\Phi_N)=c_N(0)/2$ of
\texttt{chapters/examples/cy\_d\_kappa\_stratification.tex}
Theorem~\texttt{thm:borcherds-weight-kappa-BKM-universal} is the
numerical readout of Proposition~\ref{prop:borcherds-input-normalisation};
the frontier of the conjecture at $N\in\{5,6,7,8\}$ is the
open case of whether the corresponding $\Gamma_N$-paramodular form
lifts to a \emph{diagonal-divisor} form, which the GC classification
at $t,N\leq 4$ does not settle at higher orbifold index.

\subsection*{Cycle 3: The frontier forms as two-stage-functor shadows}
\label{sec:cycle-3}

\paragraph{Attack.}
The two-stage factorisation
$\Phi_3=\Sp_{\Sigma_2,C}\circ\Phi^{\mathrm{FA}}_3$ of
\texttt{notes/wave12\_u1\_two\_stage\_functor.tex}
Definition~\texttt{def:two-stage-phi} presents the native
$E_3$-holomorphic-factorisation-algebra $\Phi^{\mathrm{FA}}_3(K3\times E)$
as the \emph{one} canonical object, with different BKMs as
specialisations along different $\Sigma_2$-cycles.
Does each frontier form fit into this picture?

\paragraph{Ghost theorem.}
Each frontier form corresponds to a specialisation
$\Sp_{\Sigma_2^{(g_N,h_M)},E}$ where $\Sigma_2^{(g_N,h_M)}\subset
X=(S\times E)/(\mathbb Z/N\mathbb Z)$ is the fibre class of the
quotient elliptic fibration compatible with $\langle g_N,h_M\rangle$.
The denominator of the resulting $E_1$-chiral shadow is the
reciprocal-square-root of the DT zeta $Z^X_{L,h_M}$
\emph{(Oberdieck--Pixton 2018)} which, by Conjecture~1 of the paper,
is the corresponding diagonal-divisor paramodular form.

\begin{conjecture}[Frontier forms as CY$_3$ specialisations]
\label{conj:frontier-as-specialisations}
For each of the eight GC forms
$\Phi_{N,L,(g_N,h_M)}$ of Table~\ref{tab:frontier-forms} and each
frontier index $N\in\{5,6,7,8\}$, there is a CY$_3$ orbifold
$X=(S\times E)/(\mathbb Z/N\mathbb Z)$ with lattice polarisation $L$
and a $(g_N,h_M)$-symplectic cover, a fibre cycle $\Sigma_2^{(g_N,h_M)}
\subset X$, and an identification
\[
\Sp_{\Sigma_2^{(g_N,h_M)},E}
\bigl(\Phi^{\mathrm{FA}}_3(X)\bigr)
\;\simeq\;
V^{(N)}_L\ \text{($E_1$-chiral algebra on }E\text{)},
\]
such that $V^{(N)}_L$ has BKM-type denominator
$\Phi_{N,L,(g_N,h_M)}=c\cdot\bigl(-Z^X_{L,h_M}\bigr)^{-1/2}$
and carries the super-structure of
$\fg_{L,(g_N,h_M)}\subset\fg_{\Delta_5}$.
\end{conjecture}

\begin{remark}[Native vs shadow at $N\geq2$]
\label{rem:native-shadow-frontier}
The slogan \emph{``many BKMs from one CY$_3$''}
(Corollary~\texttt{cor:many-BKMs} of
\texttt{notes/wave12\_u1\_two\_stage\_functor.tex}) extends to the
twisted case: a single CY$_3$ $K3\times E$ carries a canonical
$E_3$-factorisation algebra $\Phi^{\mathrm{FA}}_3(K3\times E)$, and
the eight GC forms plus the four open frontier forms at $N=5,6,7,8$
are thirteen specialisations along thirteen $(\Sigma_2,C)$-cycle
choices, with Nikulin-admissibility controlling which cycle pairs
admit a compatible orbifold structure.
The two-stage functor renders ``six routes'' (or thirteen routes)
into thirteen specialisations of the \emph{same} $E_3$-holomorphic FA,
all denominator-compatible with the Borcherds-Ray super-correction.
\end{remark}

\paragraph{Heal.}
Each frontier form is a distinct sibling BKM, not a
single-$\Phi$ multi-output; the two-stage discipline preserves
$\kcat(K3\times E)=\chi(\cO_{K3\times E})=0$ (K\"unneth on the total
space, not $\chi(\cO_{K3})=2$ on the fibre), and $\kBKM$ reads off
the weight of the corresponding paramodular form via $c_N(0)/2$.

\subsection*{Cycle 4: Frontier integrality and reflective conditions}
\label{sec:cycle-4}

\paragraph{Attack.}
A Borcherds lift converges if and only if three conditions hold:
(a) integrality of the Fourier coefficients of the Jacobi input;
(b) positivity of the real-quadratic imaginary-cone structure;
(c) reflectivity of the output lattice.
For $N=1$, (a)--(c) hold: $\phi_{0,1}$ has integral Fourier coefficients
$f(n,l)\in\mathbb Z$ (paper p.~9, line 712), the Gritsenko cone
$P_{II,\mathrm{prim}}=\{\delta_1,\delta_2,\delta_3\}$ is hyperbolic
of signature $(2,1)$ (paper \S4), and $\Lthrtwo$ is reflective
\cite[\S4]{GritsenkoNikulin1998}.
Do (a)--(c) survive at $N\geq 2$?

\paragraph{Ghost theorem.}
For $N\in\{2,3,4\}$, the twisted elliptic genera $F^{(g_N,h_M)}_L$
have integral coefficients after division by the cyclic-group
normaliser $|\langle g_N,h_M\rangle|=N$, and the corresponding
lattices $L\oplus\Lambda^{1,1}\oplus[2]$ (with
$\mathrm{rk}(L)=24-a_g^{M_{24}}$ where $a_g^{M_{24}}$ is the cycle
index at $[g_N]\in M_{24}$) are still hyperbolic/reflective after
the orbifold lattice-polarisation adjustment
\cite{Mukai1988,Chen-Oberdieck2020}.
For $N\in\{5,7\}$, the Mathieu frame shapes $1^4 5^4$ and
$1^3 7^3$ give Borcherds inputs with integral Fourier coefficients
and reflective images, but the output lies outside the GC
classification $t,N\leq 4$: these are the \emph{open} frontier
cases where the automorphic form exists and satisfies the BKM
denominator identity, but does \emph{not} have a diagonal-divisor
paramodular presentation and hence does not correspond to a
GC form.
For $N\in\{6,8\}$, the cycle-shape structure
$1^1 2^1 3^1 6^1$ and $1^1 2^1 4^1 8^1$ introduces
\emph{multiple cusps} and the Borcherds lift picks up a product over
divisor-lattice factors (see cycle~5); integrality is preserved,
but reflectivity is split across several sublattices.

\begin{proposition}[Frontier integrality and reflective closure]
\label{prop:frontier-integrality}
For every Nikulin-admissible pair $(g_N,h_M)$ with $N,M\leq 8$ and
every $(g_N,h_M)$-compatible polarisation $L$, the Fourier
coefficients $c_{F^{(g_N,h_M)}_L}(n,l)$ are integers, the
lattice $L\oplus\Lambda^{1,1}\oplus[2]$ is real-quadratic of
signature $(\mathrm{rk}(L)+2,2)$, and the signature-$(2,0)$ Weyl cone
admits a Weyl vector $\rho_{(g_N,h_M)}$ satisfying
$(\rho,\delta_i^{(N)})=-1$ on the Nikulin-reduced simple-root set
$\{\delta_i^{(N)}\}\subset L\oplus\Lambda^{1,1}\oplus[2]$.
Consequently the Borcherds product
$\Phi_{F^{(g_N,h_M)}_L}$ converges on the tube domain of the
complexified positive cone and defines an automorphic form on
$\OO(L\oplus\Lambda^{1,1}\oplus[2])_+$ of the weight
$c_{F^{(g_N,h_M)}_L}(0,0)/2$.
\end{proposition}

\begin{proof}[Proof sketch]
Integrality: the coefficients of $F^{(g_N,h_M)}_L$ in the
weak-Jacobi basis $\{\phi_{0,1},\phi_{-2,1}E_2\}$ are $M_{24}$-class
characters, hence algebraic integers in the $N$-th cyclotomic field;
Fourier coefficients are further reduced to $\mathbb Z$ by the
ring-of-integers normalisation \cite[Lem.~4.1]{GaberdielEtAl2012}.
Real-quadratic structure: the orbifold lattice-polarisation
$L\subset\mathrm{NS}(S)^{\langle g_N,h_M\rangle}$ inherits
signature $(1,\mathrm{rk}(L)-1)$ from the Picard structure, and
summing with $\Lambda^{1,1}\oplus[2]$ yields total signature
$(\mathrm{rk}(L)+2,2)$.
Weyl-vector existence: Proposition~A3a of
\texttt{notes/wave13\_a3\_gkm\_presentation\_etingof.tex} shows
$\rho=\tfrac12(\delta_1+\delta_2+\delta_3)$ at $N=1$; the
$(g_N,h_M)$-invariance of the Gram matrix on
$\{\delta_i^{(N)}\}=\langle g_N,h_M\rangle\cdot\{\delta_1,\delta_2,\delta_3\}$
secures the analogue for each $N$ admitting a Nikulin-reduction.
Convergence: Kac's asymptotic estimate
\cite[\S11]{Kac1990} together with the integrality and
real-quadraticity gives $f(n,l)\sim O(\exp(\sqrt{4n-l^2}))$, hence
Borcherds-product convergence on any neighbourhood of the
zero-dimensional cusp.
\end{proof}

\paragraph{Heal.}
Frontier integrality is uniformly satisfied for
$N\in\{1,2,3,4,5,7\}$; at $N\in\{6,8\}$ the divisor-factorisation
of the cycle shape forces a decomposition of the Borcherds product
over the divisor-lattice tower, not a failure of integrality
itself.
This is the mathematical source of Conjecture~1's stopping at
$N\leq 4$: not that integrality fails, but that the GC
diagonal-divisor classification only captures the single-cusp
Borcherds lifts and misses the multi-cusp siblings at
$N\in\{5,6,7,8\}$.

\subsection*{Cycle 5: The $N\in\{6,8\}$ divisor-tower structure}
\label{sec:cycle-5}

\paragraph{Attack.}
Slides 62--63 repeat the phrase \emph{``Case $N=6$ For all
$s\in\{0,1,2,3\}$''} twice (lines 947 and 949), and slide 64 reads
\emph{``Case $N=8$''} twice (lines 952--953).
The doubling of the slide is not an artefact of the transcription;
it mirrors the two \emph{independent} orbifold frameshapes at each
composite $N$: at $N=6$ the frame shapes $1\,2\,3\,6$ and $6^4$ give
distinct Mathieu classes, and at $N=8$ the frame shapes $1\,2\,4\,8$,
$2^2 4^2$, $8\,4\,2\,1$ give distinct classes.

\paragraph{Ghost theorem.}
Each composite-$N$ case decomposes as a product over the divisor-lattice
factors of the cycle shape, reflecting the product decomposition
of the $\eta$-function $\eta_{\text{frame}}(\tau)=\prod_i\eta(k_i\tau)^{a_i}$
in the Mathieu frame expansion.
The Borcherds lift distributes over this product as
\[
\Phi_{F^{(g_N,h_M)}_L}(Z)
=\prod_{k|N}
\Phi_{F^{(g_k,h_k)}_{L_k}}(Z|_{\text{cusp}_k})
\cdot
\Phi_{\mathrm{mix}_{(g_N,h_M)}}(Z).
\]
At $N=6$, divisors $\{1,2,3,6\}$; at $N=8$, divisors $\{1,2,4,8\}$.

\begin{proposition}[Divisor-tower decomposition at $N\in\{6,8\}$]
\label{prop:divisor-tower}
For $N\in\{6,8\}$, the frontier form $\Phi_N$ factors as a product
of divisor-indexed paramodular forms $\Phi_{k,(g_k,h_k)}$ for $k|N$
together with a mixing Borcherds-product term
$\Phi_{\mathrm{mix}_{(g_N,h_M)}}$ obtained from the non-trivial
$M_{24}$-commutator structure of $\langle g_N,h_M\rangle$.
The weight decomposes additively:
$\kBKM(\Phi_N)=\sum_{k|N}\kBKM(\Phi_{k,(g_k,h_k)})
+\kBKM(\Phi_{\mathrm{mix}})$,
with each summand being $c_k(0,0)/2$ for the corresponding
Jacobi input.
\end{proposition}

\paragraph{Heal.}
The doubled slide at $N=6$ and $N=8$ witnesses that composite-$N$
cases require two independent Borcherds lifts (one per cycle-shape
class) rather than one.
This is the source of the \emph{open} status at $N\in\{5,6,7,8\}$
of Table~\ref{tab:frontier-forms}: unlike the clean $N\in\{1,2,3,4\}$
cases where one Borcherds input yields one GC form, the composite
cases yield an unfactorised family of Siegel forms outside the GC
classification, and Conjecture~1 does not yet assert a canonical
preferred representative.

\subsection*{Connection to the two-stage functor and
$\kBKM=c_N(0)/2$ universality}
\label{sec:connection}

The frontier forms all fit into the specialisation lattice of the
native $E_3$ holomorphic factorisation algebra:
\[
\Phi^{\mathrm{FA}}_3(K3\times E)
\;\xrightarrow{\;\Sp_{\Sigma_2^{(g_N,h_M)},E}\;}\;
V^{(N)}_L
\;\xrightarrow{\;\text{denominator}\;}\;
\Phi_{N,L,(g_N,h_M)}
\;\in\;M_{c_N(0)/2}(\Gamma_N,\nu_N).
\]
The frontier extension $N\in\{5,6,7,8\}$ is the claim that this
specialisation lattice is larger than the GC
$N\leq 4$ table: the two-stage functor's cycle moduli
$\mathrm{Cyc}_2(K3\times E)\times\mathrm{Curves}(K3\times E)$ admits
Nikulin-compatible cycles at every $N\leq 8$, and the
corresponding shadow BKM algebras
$\fg_{L,(g_N,h_M)}\subset\fg_{\Delta_5}$ form a lattice of
sub-superalgebras of $\fg_{\Delta_5}$ with Mathieu-moonshine
root-multiplicity patterns.

Three consequences:
\begin{enumerate}[label={\rm(C\arabic*)}]
\item The CLAUDE.md key fact
``$\kBKM=\kch+\chi(\cO_{\mathrm{fiber}})$ is FALSE at every
$N\in\{1,2,3,4,6\}$'' extends to $N\in\{5,7,8\}$: the native-shadow
distinction kills the additive-fibre identity uniformly;
$\kBKM(\Phi_N)=c_N(0)/2$ is a Stage~2 readout of the cycle
$\Sigma_2^{(g_N,h_M)}$, while $\chi(\cO_E)=0$ and $\kch(K3\times E)=0$
are Stage~1 $\kappa_{\mathrm{cat}}$-readouts of the native FA.
\item The ``six routes to $G(K3\times E)$'' is twelve routes in
the frontier extension: the four GC$_{\leq4}$ forms, the three GC
Hecke variants, the one GC split-pair form, and the four open
frontier forms at $N\in\{5,6,7,8\}$, each a specialisation along a
distinct cycle pair $(\Sigma_2^{(g_N,h_M)},E)$.
\item The CY-A$_3$ equivalence at $N=1$ specialises to the
$\Sp_{K3,E}$-route; its Nikulin-twist-generalisation at $N\geq 2$
is the frontier conjecture: a family of chiral bar-cobar
equivalences indexed by $(\Sigma_2^{(g_N,h_M)},E)$, each realised
as a BKM denominator identity in a distinct paramodular or
Hecke-congruence weight.
\end{enumerate}

\subsection*{Physical interpretation}
\label{sec:physical}

The twisted elliptic genera $F^{(g_N,h_M)}_L$ are symmetric-orbifold
CFT partition functions at central charge $c=6$ of the
$K3/(\mathbb Z/N\mathbb Z)$-sigma model with Mathieu $M_{24}$
symmetry; the Borcherds lift is the second-quantised string
partition function on $T^*(K3/(\mathbb Z/N\mathbb Z))\times E$, by
Dijkgraaf--Moore--Verlinde--Verlinde 1997 and the
Harvey--Moore refinement.
The frontier superalgebra $\fg_{L,(g_N,h_M)}$ is the BPS algebra of
D-brane states on the orbifold CY$_3$, graded by the K3 Mukai lattice
and twisted by the $\langle g_N,h_M\rangle$-orbifold character;
its sub-superalgebra structure
$\fg_{L,(g_N,h_M)}\subset\fg_{\Delta_5}$ reflects the embedding of the
orbifold's BPS spectrum into the untwisted $K3\times E$ spectrum.
The open frontier at $N\in\{5,6,7,8\}$ corresponds to string
compactifications with $M_{24}$-twining characters outside the
diagonal-divisor paramodular locus: physically admissible, but not
captured by the diagonal-divisor classification of
Gritsenko--Cl\'ery.

\subsection*{Claim status}
\label{sec:status}

\begin{description}
\item[\ClaimStatusTheorem] The $N=1$ anchor
(Theorem~\ref{thm:n1-weight-reconciliation}): three equivalent
presentations of $\Delta_5$, all with $\kBKM=5=c_1(0,0)/2$; proved
in Lorgat 2020 + Gritsenko--Nikulin 1995 + Borcherds 1998.
\item[\ClaimStatusDerivation] Proposition~\ref{prop:borcherds-input-normalisation}
(Borcherds-input normalisation): derivable from
Eichler--Zagier + Borcherds 1998 Thm.~13.3 + Mathieu moonshine
\cite{EOT2011}; the Frame-shape constant $c_N(0,0)$ and Hecke
correction are explicit for each $M_{24}$-class.
\item[\ClaimStatusDerivation] Proposition~\ref{prop:frontier-integrality}
(frontier integrality + reflective closure): derivable via
cyclotomic-integer Fourier analysis + Nikulin's lattice-polarisation
bound + Kac's convergence estimate.
\item[\ClaimStatusDerivation] Proposition~\ref{prop:divisor-tower}
(divisor-tower decomposition at $N\in\{6,8\}$): derivable from
the $\eta$-frame-shape product decomposition + Borcherds-product
distributivity over lattice direct sums.
\item[\ClaimStatusConjectured] Conjecture~\ref{conj:frontier-as-specialisations}
(frontier-forms-as-specialisations): conditional on Stage~1 of the
two-stage functor at $d=3$ (pinned by Kontsevich--Tamarkin
$E_3$-formality + Costello--Li hCS) and on the cycle-moduli
structure of $\mathrm{Cyc}_2(K3\times E)\times\mathrm{Curves}(K3\times E)$.
\item[\ClaimStatusConjectured] The eight GC entries at $N\in\{1,2,3,4\}$
of Table~\ref{tab:frontier-forms} (and their Hecke/split-pair variants):
conditional on Conjecture~1 of Lorgat 2020 p.~1, matched case by
case to the Oberdieck--Pixton DT zeta function.
\item[\ClaimStatusOpen] The four frontier entries at
$N\in\{5,6,7,8\}$: not realised by Gritsenko--Cl\'ery's
diagonal-divisor classification; the Borcherds lift exists by
Proposition~\ref{prop:frontier-integrality} but the identification
with a DT reciprocal-square-root requires the orbifold-DT
counterpart of Oberdieck--Pixton, which is open at the
composite-$N$ cases.
\end{description}

\subsection*{New cache pattern to register}
\label{sec:cache-pattern}

\paragraph{AP-CY-FRONT1 (frontier-form specialisation vs
diagonal-divisor classification).}
The eight GC forms are the $N\leq 4$ specialisations of
$\Phi^{\mathrm{FA}}_3(K3\times E)$ along Nikulin-admissible cycles;
the frontier at $N\in\{5,6,7,8\}$ is a natural BKM-denominator
extension, not captured by diagonal-divisor classification because
the classification filters on single-cusp paramodular structure.
\emph{Wrong claim}: ``all BKM siblings of $\fg_{\Delta_5}$ lie in
the GC diagonal-divisor table.''
\emph{Ghost theorem}: the Borcherds-lift image of twisted K3
elliptic genera forms a lattice indexed by Nikulin-admissible
$(g_N,h_M)$ pairs with $N,M\leq 8$; the GC table is the
single-cusp $N\leq 4$ section of this lattice.
\emph{Correct relationship}: GC classification $=$ $N\leq 4$
diagonal-divisor Borcherds image; frontier at $N\in\{5,6,7,8\}$ $=$
multi-cusp Borcherds image (composite-$N$ divisor-tower
decomposition).
\emph{Cross-reference}: AP-CY144 (native vs shadow), AP-CY154
(canonical FA vs specialisation shadow on $K3\times E$),
AP-CY155 (imaginary-simple super-stratification on
$\fg_{\Delta_5}$), and
\texttt{notes/wave12\_u1\_two\_stage\_functor.tex}
Corollary~\texttt{cor:many-BKMs}.

\subsection*{Summary: top five frontier forms + conjectures +
wave-14 frontier}
\label{sec:summary}

\textbf{Top five frontier forms} (ordered by completeness of the
BKM-denominator picture).
\begin{enumerate}
\item $\Delta_5\in M_5(\Sp_4(\mathbb Z),\nu_2)$,
$\kBKM=5$, BKM $\fg_{\Delta_5}$: the $N=1$ anchor; proved in
Lorgat 2020 Theorem~3; proved-three-ways reconciliation in
Theorem~\ref{thm:n1-weight-reconciliation}.
\item $M_2(\Gamma_2,\nu_4)$, $\kBKM=2$, BKM
$\fg_{L_2,(g_2,g_2)}\subset\fg_{\Delta_5}$: the $N=2$
anchor-of-frontier; Borcherds input $F^{(g_2,g_2)}_{L_2}$ with
$M_{24}$-frame shape $1^8 2^8$; conditional on Conjecture~1.
\item $M_3(\Gamma_1(2),\nu_2)$, $\kBKM=3$, Hecke-congruence
$N=2$ sibling; Borcherds input $F^{(g_2,\mathrm{id})}_{L_2}$;
paramodular-vs-Hecke contrast at fixed $N$ showcases the
Stage-2 cycle-moduli dependence of the chiral shadow.
\item $M_1(\Gamma_3,\nu_6)$, $\kBKM=1$, BKM
$\fg_{L_3,(g_3,g_3)}\subset\fg_{\Delta_5}$: the $N=3$ anchor;
lowest GC-classified weight admitting a paramodular presentation;
Borcherds input $F^{(g_3,g_3)}_{L_3}$ with frame shape $1^6 3^6$.
\item $M_{1/2}(\Gamma_4,\nu_8)$, $\kBKM=1/2$, BKM
$\fg_{L_4,(g_4,g_4)}\subset\fg_{\Delta_5}$: the unique half-integer
Gritsenko-Cl\'ery form; frame shape $2^4 4^4$; minimal paramodular
weight in the GC table.
\end{enumerate}

\textbf{Conjectures} (Lorgat 2020 Conjecture~1 + extensions).
\begin{itemize}
\item \emph{Diagonal-divisor-DT correspondence}: all eight GC forms
arise as $(-Z^X_{L,h_M})^{-1/2}$ for matched CY$_3$ data
$(S,L,g_N,h_M,E,e_0)$ with $N\leq 4$ (Lorgat 2020 Conjecture~1).
\item \emph{Frontier-BKM extension}: the specialisation
lattice of $\Phi^{\mathrm{FA}}_3(K3\times E)$ along
Nikulin-admissible cycles at $N\in\{5,6,7,8\}$ produces four
additional sibling BKM algebras $\fg_{L_N,(g_N,g_N)}$ with
Borcherds-denominator weights $c_N(0)/2$, not diagonal-divisor
but otherwise paramodular or Hecke-congruence
(Conjecture~\ref{conj:frontier-as-specialisations}).
\item \emph{Universal Borcherds weight} ($\kBKM(\Phi_N)=c_N(0)/2$):
Stage~2 invariant of the two-stage functor;
independent of cycle shape across the frontier
(\texttt{chapters/examples/cy\_d\_kappa\_stratification.tex}
Theorem~\texttt{thm:borcherds-weight-kappa-BKM-universal}).
\end{itemize}

\textbf{Wave-14 frontier} (next research moves).
\begin{itemize}
\item Identify $\Sigma_2^{(g_N,h_M)}\subset K3\times E$ explicitly
for each of the eight GC data; write down the cycle class in
$H_4(K3\times E,\mathbb Z)$ compatible with the
$\langle g_N,h_M\rangle$-orbifold.
\item Compute the twisted elliptic genus $F^{(g_N,h_M)}_L$ for
every Nikulin-admissible pair at $N\leq 4$ explicitly in the
$\phi_{0,1}\oplus\phi_{-2,1}E_2$-basis, and extract
$c_N(0,0)\in 2\mathbb Z$ as a combinatorial readout of the
frame-shape $M_{24}$-character.
\item For each of the frontier $N\in\{5,6,7,8\}$ cases, execute
Proposition~\ref{prop:frontier-integrality} numerically: compute
$F^{(g_N,g_N)}_{L_N}$, the Borcherds product
$\Phi_{F^{(g_N,g_N)}_{L_N}}$, and its weight
$c_N(0,0)/2$; match against known
paramodular/Hecke-congruence weights via Gritsenko-lift tables.
\item For $N\in\{6,8\}$, resolve the divisor-tower decomposition
of Proposition~\ref{prop:divisor-tower} explicitly: specify the
divisor-indexed restrictions $Z|_{\text{cusp}_k}$ and the mixing
term $\Phi_{\mathrm{mix}}$.
\item Physical interpretation at $N\in\{5,7\}$ (prime case):
the $M_{24}$ twining classes $[g_5]$ and $[g_7]$ are the
non-Leech Niemeier-lattice classes; their BKM siblings
$\fg_{L_5,(g_5,g_5)}$ and $\fg_{L_7,(g_7,g_7)}$ are BPS algebras
of rare stringy orbifolds with no 6D supergravity dual in the
standard CHL lattice; worth identifying in heterotic-type~II
duality.
\item Close the CY-A$_3$ twisted-equivalence conjecture for each
$N\leq 8$: $\mathrm{Sp}_{\Sigma_2^{(g_N,h_M)},E}(\Phi^{\mathrm{FA}}_3(X))$
as chiral bar--cobar equivalence with reference data from the
orbifold DT.
\end{itemize}

\subsection*{Cross-references}
\label{sec:cross-refs}

Primary sources (text): \texttt{notes/wave13\_source\_paper.txt}
(Lorgat 2020, 783 lines), \texttt{notes/wave13\_source\_slides.txt}
lines 800--953 (last 10 slides, $N$-case enumeration).
Two-stage functor:
\texttt{notes/wave12\_u1\_two\_stage\_functor.tex}
Def.~\texttt{def:two-stage-phi},
Corollary~\texttt{cor:many-BKMs}.
Kazhdan-voice BKM universality:
\texttt{notes/wave12\_a2\_k3xE\_BKM\_kazhdan.tex}
(Cycle~1 $N=1$ anchor, Cycle~2 $N\geq 2$).
Cache row F8 (native/shadow inversion):
\texttt{appendices/first\_principles\_cache.md} row F8.
Cache rows A3a--c (Weyl vector, super-stratification, correction
necessity): \texttt{appendices/first\_principles\_cache.md}
rows A3a, A3b, A3c.
Reader-facing $\kappa_\bullet$-spectrum table
(subscripts $\bullet\in\{\mathrm{ch},\mathrm{cat},\mathrm{BKM},\mathrm{fiber}\}$):
\texttt{chapters/examples/cy\_d\_kappa\_stratification.tex}.

\medskip
\noindent
\emph{Authorship: Raeez Lorgat.}
